\documentclass[11pt]{amsart}

\usepackage[T1]{fontenc}
\usepackage{lmodern}
\usepackage{microtype}
\usepackage{amsmath,amssymb,amsthm}
\usepackage[hidelinks]{hyperref}

\hypersetup{
  pdftitle={The Gurvich--Naumova GM-Strategy Conjecture Fails with Three Piles}
}

\newtheorem{theorem}{Theorem}
\newtheorem{corollary}[theorem]{Corollary}
\newcommand{\NIM}{\operatorname{NIM}}

\title[The GM-strategy conjecture]
{The Gurvich--Naumova GM-Strategy Conjecture Fails with Three Piles}
\date{}

\subjclass[2020]{91A46, 91A06, 91A10}
\keywords{exact slow NIM, GM strategy, Nash equilibrium, counterexample}

\begin{document}

\begin{abstract}
Gurvich and Naumova conjectured that, in the $\ell$-player version of
exact slow $\NIM(n,n-1)$, the profile in which every player follows the
GM strategy is a uniform Nash equilibrium. We disprove the conjecture
for every $\ell\ge 3$, already with three piles. From the initial
position $(2,\ell+1,\ell+1)$, all-GM play lasts $\ell+2$ moves, whereas
player~1 can keep the pile of size $2$ on the first move and follow GM
thereafter, producing a play of length $\ell+1$. Player~1 remains a
winner and improves his payoff by $1/(\ell-1)$. We also show that three
is the least possible number of piles for such a counterexample.
\end{abstract}

\maketitle

\section{Counterexample and minimality}

Gurvich and Naumova conjectured that the $\ell$ GM strategies form a
uniform Nash equilibrium in the $\ell$-player game
$\NIM(n,n-1)$~\cite[Section~6.5, Conjecture~1]{GN}, whose statement
reads: ``The set of $\ell$ GM-strategies form a (uniform) Nash
equilibrium.'' (Here and below the numbering is that of the e-print
arXiv:2312.08382.) They note that the conjecture holds for $\ell=2$.
We disprove it for every $\ell\ge 3$, already for $n=3$.

A position of $\NIM(3,2)$ is a nonnegative triple, written in
nondecreasing order. A move keeps one coordinate and decreases the
other two by one; the decreased coordinates must therefore be
positive. In $\NIM(n,n-1)$ a position is terminal when at least two of
its coordinates are non-positive~\cite[Section~6.5]{GN}; for $n=3$ this
says that at most one coordinate is positive, and for $n=2$ that the
only terminal position is $(0,0)$. Players $1,\ldots,\ell$ move
cyclically, with player~1 first, and the player whose turn follows the
last move is the loser. For a
fixed constant $C$ larger than the length of every play from the given
initial position, the loser has payoff $L-C$, while each winner has
payoff
\[
  \frac{C-L}{\ell-1},
\]
where $L$ is the length of the play.

For $\NIM(3,2)$, the GM strategy keeps a smallest coordinate divisible
by $\ell$, ties being broken in favour of the rightmost such coordinate
in the nondecreasing ordering; if no coordinate is divisible by $\ell$,
it keeps a largest coordinate.

\begin{theorem}\label{thm:main}
For every integer $\ell\ge 3$, in the $\ell$-player game
$\NIM(3,2)$ initialized at
\[
  x^0=(2,\ell+1,\ell+1),
\]
the all-GM profile is not a Nash equilibrium. Player~1 has a
unilateral deviation that increases his payoff by $1/(\ell-1)$.
\end{theorem}

\begin{proof}
At $x^0$, no coordinate is divisible by $\ell$, so GM keeps a largest
coordinate. At the next position, the coordinate $\ell$ is the unique
multiple of $\ell$ and is kept. Thereafter the zero coordinate must be
kept. Hence all-GM play is
\[
\begin{aligned}
(2,\ell+1,\ell+1)
&\longrightarrow (1,\ell,\ell+1)
 \longrightarrow (0,\ell,\ell)\\
&\longrightarrow (0,\ell-1,\ell-1)
 \longrightarrow \cdots
 \longrightarrow (0,0,0).
\end{aligned}
\]
The tail from $(0,\ell,\ell)$ has $\ell$ moves, so the total length is
$L_{\mathrm{GM}}=\ell+2$. Thus player~3 is the loser and player~1 is a
winner.

Let $\sigma^{\mathrm{GM}}$ denote the GM strategy. Define a strategy
$\tau_1$ for player~1 that agrees with $\sigma^{\mathrm{GM}}$ at every
position except $x^0$, where it keeps the pile of size $2$. Since
$\ell\ge 3$, the coordinate $2$ is not a multiple of $\ell$, so at
$(2,\ell,\ell)$ GM keeps a coordinate $\ell$; at $(1,\ell-1,\ell)$ the
coordinate $\ell$ is the unique multiple of $\ell$ and is kept;
thereafter the zero coordinate must be kept. Hence, against GM play by
all other players, the resulting path is
\[
\begin{aligned}
(2,\ell+1,\ell+1)
&\longrightarrow (2,\ell,\ell)
 \longrightarrow (1,\ell-1,\ell)\\
&\longrightarrow (0,\ell-2,\ell)
 \longrightarrow \cdots
 \longrightarrow (0,0,2).
\end{aligned}
\]
The tail from $(0,\ell-2,\ell)$ has $\ell-2$ moves, so
$L_{\mathrm{dev}}=\ell+1$. Now player~2 is the loser, and player~1 is
again a winner. His payoff therefore changes by
\[
  \frac{C-(\ell+1)}{\ell-1}
  -\frac{C-(\ell+2)}{\ell-1}
  =\frac{1}{\ell-1}>0.
\]
Both ties encountered above, at $x^0$ and at $(2,\ell,\ell)$, are
between coordinates of equal value, so neither of the two plays depends
on the tie-breaking rule. Thus $\tau_1$ is a profitable unilateral
deviation.
\end{proof}

The strategies in question do not depend on the initial position: ``By
definition, the GM-strategies are uniform, that is, independent of
$x^0$''~\cite[Section~6.5]{GN}. A Nash equilibrium in GM-strategies
would therefore have to hold at every initial position, so the failure
at the single position $x^0=(2,\ell+1,\ell+1)$ refutes Conjecture~1 for
every $\ell\ge 3$.

The hypothesis $\ell\ge 3$ is used only to ensure that $2$ is not a
multiple of $\ell$ (and that player~3 exists). For $\ell=2$ the GM move
at $x^0$ already keeps the pile of size $2$, so $\tau_1$ coincides with
the GM strategy and the construction degenerates.

\begin{corollary}\label{cor:minimal}
Among the games $\NIM(n,n-1)$ with $n\ge 2$, three is the least number
of piles for which the all-GM profile can fail to be a Nash equilibrium.
\end{corollary}

\begin{proof}
For $n=2$, the game is $\NIM(2,1)$. Every move decreases the total
number of stones by exactly one, and by the terminality convention
above the only terminal position is $(0,0)$. Thus
every play from $(a,b)$ has length $a+b$, independently of all
strategic choices. The loser and all payoffs are consequently fixed,
so no unilateral deviation is profitable. Theorem~\ref{thm:main}
shows that $n=3$ is attainable.
\end{proof}

\begin{thebibliography}{1}

\bibitem{GN}
V.~Gurvich and M.~Naumova,
\emph{Screw discrete dynamical systems and their applications to exact
slow NIM},
Discrete Appl. Math. \textbf{358} (2024), 382--394;
e-print \href{https://arxiv.org/abs/2312.08382}{arXiv:2312.08382}.
\href{https://doi.org/10.1016/j.dam.2024.07.018}
{doi:10.1016/j.dam.2024.07.018}.

\end{thebibliography}

\end{document}